The goal of TTS synthesis is to render  naturally sounding speech signals given a text to be synthesized.
Human speech production process first translates a text (or concept) into movements of muscles associated with
articulators and speech production-related organs.
Then using air-flow from lung, vocal source excitation signals, which contain both periodic
(by vocal cord vibration) and aperiodic (by turbulent noise) components, are generated.
By filtering the vocal source excitation signals by time-varying vocal tract transfer functions
controlled by the articulators, their frequency characteristics are modulated.
Finally, the generated speech signals are emitted.
The aim of TTS is to mimic this process by computers in some way.

TTS can be viewed as a sequence-to-sequence mapping problem;
from a sequence of 
discrete symbols (text) to a real-valued time series (speech signals).
A typical TTS pipeline has two parts; 1) text analysis and 2) speech synthesis.
The text analysis part typically includes a number of natural language processing (NLP) steps,
such as sentence segmentation, word segmentation, text normalization,
part-of-speech (POS) tagging,
and grapheme-to-phoneme (G2P) conversion.
It takes a word sequence as input and outputs
a phoneme sequence with a variety of linguistic contexts.
The speech synthesis part takes the context-dependent phoneme sequence as its input and outputs a synthesized speech waveform.
This part typically includes prosody prediction and speech waveform generation.

\begin{figure}[!t]
  \centering
  \includegraphics[width=0.6\textwidth]{Figures/hts_flow.pdf}
  \caption{Outline of statistical parametric speech synthesis.}
  \label{fig:outline_spss}
\end{figure}%mm

There are two main approaches to realize the speech synthesis part; non-parametric, example-based approach known as concatenative speech synthesis \citep{PSOLA,ATR_nutalk,Hunt_UnitSelection_ICASSP}, and parametric, model-based approach known as statistical parametric speech synthesis \citep{yoshimura_PhD,Zen_SPSS_SPECOM}.
The concatenative approach builds up the utterance from units of recorded speech, whereas the statistical parametric approach uses a generative model to synthesize the speech.
The statistical parametric approach first extracts a sequence of vocoder parameters \citep{Vocoder} $\bo = \{ \bo_1, \dots, \bo_N \}$ from speech signals $\bx = \{ x_1, \dots, x_T \}$ and linguistic features $\bl$ from the text $W$, where $N$ and $T$ correspond to the numbers of vocoder parameter vectors and speech signals.
Typically a vocoder parameter vector $\bo_n$ is extracted at every 5 milliseconds.
It often includes cepstra \citep{UELS} or line spectral pairs \citep{LSP}, which represent 
vocal tract transfer function, and fundamental
frequency ($F_0$) and aperiodicity \citep{Kawahara_STRAIGHT_Excitation}, which represent characteristics of vocal source excitation signals.
Then a set of generative models, such as hidden Markov models (HMMs) \citep{yoshimura_PhD}, feed-forward neural networks \citep{Zen_DNN_ICASSP}, and recurrent neural networks \citep{Robinson_RNNTTS,Karaani_RTDNNTTS,Fan_BLSTM_Interspeech14}, is trained from the extracted vocoder parameters and linguistic features as
\begin{equation}
  \hat{\Lambda} = \argmax_{\Lambda} p\left( \bo \mid \bl, \Lambda \right),
\end{equation}
where $\Lambda$ denotes the set of parameters of the generative model.
At the synthesis stage, the most probable vocoder parameters are generated given linguistic features extracted from a text to be synthesized as
\begin{equation}
  \hat{\bo} = \argmax_{\bo} p ( \bo \mid \bl, \hat{\Lambda} ).
\end{equation}
Then a speech waveform is reconstructed from $\hat{\bo}$ using a vocoder.
The statistical parametric approach offers various advantages over the
concatenative one such as small footprint
and flexibility to change its voice characteristics.
However, its subjective naturalness is often significantly worse than that of the concatenative approach; synthesized speech often sounds muffled and has artifacts.
\cite{Zen_SPSS_SPECOM} reported three major factors that can degrade the subjective naturalness; quality of vocoders, accuracy of generative models, and effect of oversmoothing.
The first factor causes the artifacts and the second and third factors lead to the muffleness in the synthesized speech.
There have been a number of attempts to address these issues individually, such as developing high-quality vocoders \citep{Kawahara_STRAIGHT,Vocaine,WORLD}, improving the accuracy of generative models \citep{Zen_trjHMM_CSL,Zen_DNN_ICASSP,Fan_BLSTM_Interspeech14,Uria_TrajectoryRNADE_ICASSP2015}, and compensating the oversmoothing effect \citep{Toda_MLGV_IEICE,Takamichi_ModulationSpectrum_TASLP}.
\cite{Zen_LSTMprod_Interspeech} showed that state-of-the-art statistical parametric speech syntheziers matched state-of-the-art concatenative ones in some languages.
However, its vocoded sound quality is still a major issue.

Extracting vocoder parameters can be viewed as estimation of a generative model parameters given speech signals \citep{LPC,UELS}.
For example, linear predictive analysis \citep{LPC}, which has been used in speech coding, assumes that the generative model of speech signals is 
a linear auto-regressive (AR) zero-mean Gaussian process;
\begin{align}
  x_t &= \sum_{p=1}^P a_p x_{t-p} + \epsilon_t \\
  \epsilon_t &\sim \Gauss(0, G^2)
\end{align}
where $a_p$ is a $p$-th order linear predictive coefficient (LPC) and $G^2$ 
is a variance of modeling error.
These parameters are estimated based on the maximum likelihood (ML) criterion.
In this sense, the training part of the statistical parametric approach can be viewed as a two-step optimization and sub-optimal: extract vocoder parameters by fitting a generative model of speech signals then model trajectories of the extracted vocoder parameters by a separate generative model for time series \citep{Tokuda_ASRU2011}.
There have been attempts to integrate these two steps into a single one
\citep{STAVOCO,MGELSD,Maia_WaveformModel_SSW,Kazuhiro_McepHMM_IEICE,Black_AutoEncoder_arXiv,Tokuda_CepLSTM_ICASSP2015,Tokuda_MixCepLSTM_ICASSP2016,Takaki_FFTDNN_ICASSP}.
For example, \cite{Tokuda_MixCepLSTM_ICASSP2016} integrated non-stationary, nonzero-mean Gaussian process generative model of speech signals and LSTM-RNN-based sequence generative model to a single one and jointly optimized them by back-propagation.
Although they showed that this model could approximate natural speech signals, its segmental naturalness was significantly worse than the non-integrated model due to over-generalization and over-estimation of noise components in speech signals.

The conventional generative models of raw audio signals have a number of assumptions which are inspired from the speech production, such as
\begin{itemize}
    \item Use of fixed-length analysis window; They are typically based on a stationary stochastic process \citep{LPC,UELS,Poritz_ARHMM_ICASSP82,Juang_MARHMM_ICASSP85,Kameoka_MultiKernelLPC_ASJ}.  To model time-varying speech signals by a stationary stochastic process, parameters of these generative models are estimated within a fixed-length, overlapping and shifting analysis window (typically its length is 20 to 30 milliseconds, and shift is 5 to 10 milliseconds).
    However, some phones such as stops are time-limited by less than 20 milliseconds \citep{Rabiner_ASR}.
    Therefore, using such fixed-size analysis window has limitations.
    \item Linear filter; These generative models are typically realized as a linear time-invariant filter  \citep{LPC,UELS,Poritz_ARHMM_ICASSP82,Juang_MARHMM_ICASSP85,Kameoka_MultiKernelLPC_ASJ} within a windowed frame.  However, the relationship between successive audio samples can be highly non-linear.
    \item Gaussian process assumption; The conventional generative models are based on Gaussian process \citep{LPC,UELS,Poritz_ARHMM_ICASSP82,Juang_MARHMM_ICASSP85,Kameoka_MultiKernelLPC_ASJ,Tokuda_CepLSTM_ICASSP2015,Tokuda_MixCepLSTM_ICASSP2016}.  From the source-filter model of speech production \citep{Chiba_SourceFilter,Fant_SourceFilter} point of view, this is equivalent to assuming that a vocal source excitation signal is a sample from a Gaussian distribution \citep{LPC,UELS,Poritz_ARHMM_ICASSP82,Juang_MARHMM_ICASSP85,Tokuda_CepLSTM_ICASSP2015,Kameoka_MultiKernelLPC_ASJ,Tokuda_MixCepLSTM_ICASSP2016}. 
    Together with the linear assumption above, it results in assuming that speech signals are normally distributed.
    However, distributions of real speech signals can be significantly different from Gaussian.
\end{itemize}
Although these assumptions are convenient, samples from these generative models tend to be noisy and lose important details to make these audio signals sounding natural.

WaveNet, which was described in Section~\ref{sec:wavenet}, has none of the above-mentioned assumptions.
It incorporates almost no prior knowledge about audio signals, except the choice of the receptive field and $\mu$-law encoding of the signal.
It can also be viewed as a non-linear causal filter for quantized signals.
Although such non-linear filter can represent complicated signals while preserving the details, designing such filters is usually difficult \citep{NonlinearFilterDesign}.
WaveNets give a way to train them from data.

